% ═══════════════════════════════════════════════════════════════════════════
% QR-CODE: Lernpfad Kräfte zusammensetzen (UE 2)
% Physik Klasse 7 – Beamer-Projektion
% ═══════════════════════════════════════════════════════════════════════════

\documentclass[a4paper,11pt]{article}

\usepackage[utf8]{inputenc}
\usepackage[T1]{fontenc}
\usepackage{helvet}
\renewcommand{\familydefault}{\sfdefault}

\usepackage[margin=20mm]{geometry}
\usepackage{xcolor}
\usepackage{qrcode}

% Farben
\definecolor{physik}{HTML}{2c5aa0}

% Konfiguration
\newcommand{\lptitel}{Kräfte zusammensetzen}
\newcommand{\lpurl}{https://devbox42.github.io/sciencesim/physik/kl7/kraft/LP-02-kraft-groesse.html}
\newcommand{\lpurlkurz}{devbox42.github.io/.../LP-02-kraft-groesse.html}

\pagestyle{empty}

\begin{document}

\begin{center}

% Titel
\vspace*{10mm}
{\fontsize{28}{34}\selectfont\bfseries\color{physik}Lernpfad: \lptitel}

\vspace{15mm}

% QR-Code mit Rahmen
\fcolorbox{physik}{white}{\qrcode[height=100mm]{\lpurl}}

\vspace{12mm}

% URL
{\large\color{gray}\lpurlkurz}

\vspace{8mm}

% Hinweis
{\Large\color{gray}Scanne den QR-Code mit deinem Gerät}

\vspace{10mm}

% Info-Box
\fcolorbox{physik}{physik!10}{%
    \parbox{0.7\textwidth}{%
        \centering
        \vspace{3mm}
        {\large Arbeite die Schritte durch und nutze das Arbeitsblatt.}\\[2mm]
        Am Ende: Minitest nicht vergessen!
        \vspace{3mm}
    }%
}

\end{center}

\end{document}
